\section{Joint arithmetic localization of optimized block gain (NS-87)}
\paragraph{Scope.}
This is a localization theorem and a certified finite descent prescription,
not a new arithmetic lower estimate. The Hilbert-space inequality is generic;
the finite observations below use the original arithmetic atoms. The
polynomial-cutoff conclusion uses the already established NS-81 coefficient
bound and NS-83 critical-zero lower bound. It supplies neither an effective
asymptotic onset nor decay of the optimized error.

Use the fixed $q=2$ setting of NS-86 in reciprocal coordinates,
$\mathcal H=L^2((0,\infty),dt/t^2)$, with target
$\tau(t)=\boldsymbol1_{t\ge1}\log t$ and atoms
$a_n(t)=A(t/n)$, where $A(z)=\int_0^z\{u\}\,du/u$.
Let $V_N=\operatorname{span}(a_1,\ldots,a_N)$, $\Pi_N$ its orthogonal
projection, $r_N=(I-\Pi_N)\tau$, and $E_N=\|r_N\|^2$.
For $M>N$, set $\phi_n=(I-\Pi_N)a_n$, $N<n\le M$, with positive definite
Gram $H=(\langle\phi_i,\phi_j\rangle)$. Let
$h=(\langle r_N,\phi_n\rangle)$ and $\Delta=E_N-E_M=h^TH^{-1}h$.
All transposes are real; the Hermitian version is identical.

\subsection{Finite arithmetic observations}
For any $T\ge1$, put
\[
 P_T(t)=\int_t^T\log(u/t)r_N(u)\frac{du}{u^2}\quad(0<t\le T),
 \qquad P_T(t)=0\quad(t>T).
\]
These are truncated potentials, not the full $P_N$ of NS-86. Then
\begin{equation}
 q_n(T):=\langle\boldsymbol1_{t\le T}r_N,a_n\rangle
 =\frac1n\int_0^T P_T(t)\,dt-\sum_{1\le k\le T/n}P_T(kn).
 \label{eq:ns87-finite-quadrature}
\end{equation}
To prove it, use $(tP_T')'=r_N/t^2$ and
$(ta_n')'=1/n-\sum_{k\ge1}\delta_{kn}$. Integrating twice by parts
has zero upper boundary terms, since $P_T(T)=P_T'(T)=0$; the lower
ones vanish as $O(t(1+\log^2t))$. This also treats nonaligned $T$;
a summand at $kn=T$ is zero. No infinite sum or divergent integral is
separated. Equivalently,
\[
 \int_0^T P_T(t)\,dt=\int_0^T r_N(u)\,du/u
\]
by Fubini, since $\int_0^u\log(u/t)\,dt=u$.
If $D$ is the coefficient matrix of the projected atoms, the finite
observation vector is $\widehat h_T=D^T(q_1(T),\ldots,q_M(T))^T$.
One must keep the old-coordinate rows of $D$: the truncated residual is
not orthogonal to $V_N$, even though the full residual is.

\subsection{Coupled tail, complete gain and an explicit safe step}
Let
\[
 e_T=\int_T^\infty r_N(t)^2\frac{dt}{t^2},\qquad
 K_T=\left(\int_T^\infty\phi_i(t)\phi_j(t)\frac{dt}{t^2}\right),
 \quad b_T=\min\{1,\operatorname{tr}(H^{-1}K_T)\}.
\]
The trace is nonnegative. Any certified scalar upper bounds
$\overline e_T\ge e_T$, $\overline b_T\ge b_T$ can replace these quantities;
in particular $\overline b_T=1$ always works. Write $B_T=\overline b_T\overline e_T$
and $A_T=\widehat h_T^TH^{-1}\widehat h_T$. Then
\begin{align}
 (h-\widehat h_T)^TH^{-1}(h-\widehat h_T)&\le B_T,
 \label{eq:ns87-tail-metric}\\
 (\sqrt{A_T}-\sqrt{B_T})_+^2
 &\le\Delta\le(\sqrt{A_T}+\sqrt{B_T})^2.
 \label{eq:ns87-gain-interval}
\end{align}
Proof: the synthesis $U:x\mapsto\sum x_n\phi_n$ obeys $U^*U=H$,
so $Z=UH^{-1/2}$ is an isometry. If $Q_T$ is multiplication by
$\boldsymbol1_{t>T}$, the omitted vector is $U^*Q_Tr_N$. Its squared
$H^{-1}$ norm is at most $\|Z^*Q_T\|^2e_T$.
The squared operator norm is the largest eigenvalue of
$H^{-1/2}K_TH^{-1/2}$, bounded by both one and its trace. The ordinary
triangle and reverse triangle inequalities then give the gain interval.
This joint argument avoids replacing the tail by unrelated entrywise errors.
It still uses the complete arithmetic Gram; it does not construct its fast inverse.

Moreover the lower bound is constructive. For $A_T>B_T$, set
\[
 v_T=H^{-1}\widehat h_T,\quad
 W_T=\sum(v_T)_n\phi_n,\quad
 \theta_T=1-\sqrt{B_T/A_T}.
\]
The corrected residual $r_N-\theta_TW_T$ is admissible at size $M$ and
\begin{equation}
 E_N-\|r_N-\theta_TW_T\|^2
 \ge(\sqrt{A_T}-\sqrt{B_T})^2.
 \label{eq:ns87-safe-step}
\end{equation}
Indeed $\|W_T\|^2=A_T$ and
$\langle r_N,W_T\rangle\ge A_T-\sqrt{A_TB_T}$.
If $A_T\le B_T$, take the zero step. An exact dyadic
$0\le\theta\le1-\sqrt{B_T/A_T}$ guarantees at least $\theta^2A_T$;
this is the outward-rounded prescription in the verifier.
The complete new coefficients include the old-coordinate correction $Dv_T$.
No claim of unit-step descent or of preserving NS-86's convolution identity is made.

\subsection{The observation tail is negligible at a polynomial cutoff}
This is an asymptotic localization result, not a convergence theorem.
Let $\mathsf H_N=\sum_{n\le N}1/n$ and
$C_N=48\mathsf H_NN^3(N+2)$ from NS-81. If
$r_N=\tau-\sum c_na_n$, orthogonal projection and NS-81 give
\[
 \|c\|_{\ell^2}^2\le16C_N\|\Pi_N\tau\|^2\le32C_N.
\]
For every $T\ge N$, define
\[
 a=\tfrac12\log(2\pi T)+\frac{N}{6T},\qquad
 F_N(T)=\frac{2((\log T)^2+2\log T+2)
                 +64NC_N(a^2+a+\tfrac12)}{T}.
\]
Then the full residual tail satisfies $e_T\le F_N(T)$.
Indeed, the complete Stirling bound gives
$0\le A(t/n)\le a+\tfrac12\log(t/T)$ for $t\ge T$ and $n\le N$.
Cauchy--Schwarz over coefficients gives
$(\sum c_na_n(t))^2\le32NC_N(a+\tfrac12\log(t/T))^2$.
Apply $(x-y)^2\le2x^2+2y^2$ and integrate. The target integral is
$((\log T)^2+2\log T+2)/T$ and the latter integral is
$(a^2+a+1/2)/T$, proving the stated bound with every tail term retained.

Fix any $\varepsilon>0$ and put $T_N=\lceil N^{5+\varepsilon}\rceil$.
Then $F_N(T_N)=O_\varepsilon(N^{-\varepsilon}\log^3N)$.
NS-83 gives $E_N\ge c/\log N$ for all sufficiently large $N$, for some
$c>0$, unconditionally from known critical-line zeros. Thus
\begin{equation}
 \frac{F_N(T_N)}{E_N}
 =O_\varepsilon(N^{-\varepsilon}\log^4N).
 \label{eq:ns87-relative-tail}
\end{equation}
This use of NS-83 supplies no explicit onset. We never apply its
asymptotic lower bound at one of the tested finite sizes.

For any $M=M(N)>N$, take $B_T=F_N(T_N)$ in the preceding theorem.
Write $\gamma_N=(E_N-E_M)/E_N\le1$ and
$L_N=(\sqrt{A_{T_N}}-\sqrt{F_N(T_N)})_+^2$.
Putting $\epsilon_N=\sqrt{F_N(T_N)/E_N}$ yields
\[
 \left|\frac{A_{T_N}}{E_N}-\gamma_N\right|
 \le2\epsilon_N+\epsilon_N^2,\qquad
 0\le\gamma_N-\frac{L_N}{E_N}\le4\epsilon_N.
\]
For the first inequality use
$|\sqrt{A/E}-\sqrt\gamma|\le\epsilon_N$ and $\sqrt\gamma\le1$.
For the second, the square roots of $\gamma$ and $L/E$ differ by at most
$2\epsilon_N$, and both lie in $[0,1]$.
Consequently
\begin{equation}
 \log N\left(\frac{A_{T_N}}{E_N}-\gamma_N\right)\longrightarrow0,
 \qquad
 \log N\left(\gamma_N-\frac{L_N}{E_N}\right)\longrightarrow0.
 \label{eq:ns87-log-scale-localization}
\end{equation}
The simple choice $T_N=N^6$ is sufficient. No matrices or sums at that
asymptotic cutoff were assembled for this assertion: it is an analytic bound.
No sufficiency claim for $T=N^2$ follows.

\paragraph{The remaining arithmetic target.}
For $N=2^j$, $M=2N$, it is now enough to establish, for the actual optimized
coefficients, $\liminf_{j\to\infty}jA_{N^6}/E_N>0$.
The localization proves the same positive lower liminf for $j\gamma_N$,
so the product of successive error factors vanishes. This is an open sufficient
estimate, not a proved estimate or a newly discovered RH criterion.
The corresponding positive liminf need not be necessary for convergence.
NS-83 excludes eventual $\gamma_N\ge a/j$ with $a>1$; it does not
exclude $0<a\le1$ or weaker nonsummable lower scales. Equations above transfer
such leading-scale statements without a hidden infinite observation tail.
The Gram, optimized coefficients, and growing signed sums remain arithmetic
objects requiring an independent bound. The generic Hilbert argument alone
cannot distinguish the NS-74 altered-family control.

\subsection{Finite validation}
The verifier evaluates $N=8,16,32$, $M=2N$, at $T=M,4M,16M,64M$.
It uses complete Gram entries with their analytic Bernoulli remainder,
interval old-space solves, exact physical cells and a second finite-potential
calculation. No midpoint solve is used. Each cell is a combination of
$t$, $\log t$ and $1$, and the exterior $0<t<1$ is retained.
For every case the full physical tail $e_T=E_N-\int_0^Tr_N^2dt/t^2$
is enclosed via the complete Gram and separately checked against the
Stirling tail interval. At $N=32,T=4096$, the safe step guarantees more
than $95.98\%$ of full available gain; this is more than $18.15\%$ of
that finite old error, not a uniform fraction at all sizes.
At $N=32,T=64$, $A_T$ exceeds full gain by more than $15.99\%$;
discarding the tail would therefore overstate the guaranteed reduction.
Two arithmetic precisions and a doubled analytic Gram cutoff replay the same
finite problems. Changing the physical observation cutoff is a different
localized observation problem and is not called a precision replay.
